
The sequence \((*)\) is an exact sequence of \(H\)-groups,
provided that \(H\) acts on \(V^{(p)}\) through the linear
representation
\[
  H\xrightarrow{v}\operatorname{GL}(V)
    \xrightarrow{F}\operatorname{GL}(V^{(p)}).
\]
Since \(k\) is perfect, the morphism \(F:V\to V^{(p)}\) induces
a surjective map
\[
  V(k)\longrightarrow V^{(p)}(k).
\]
It follows from Appendix I, \SGARefLink{sga3:i:subsection:2-1}{2.1} that the exact sequence \((*)\)
defines an exact sequence
\[
  \tag{**}\label{sga3:xvii:proposition:5-4-1:equation:double-star}
  H^1(H,V^{(p)})\longrightarrow
  \operatorname{Ext}_{\mathrm{alg}}(H,{}_F V)
  \longrightarrow
  \operatorname{Ext}_{\mathrm{alg}}(H,V).
\]

We shall show that
\[
  \operatorname{Ext}_{\mathrm{alg}}(H,V)=0.
\]
Let \(E'\) be an algebraic group which is an extension of \(H\)
by \(V\):
\[
  1\longrightarrow V\longrightarrow E'
    \xrightarrow{h}H\longrightarrow1.
\]
The scheme \(E'\) is a torsor over \(H\) under
\(\mathbb G_a^r\), and hence defines an element of
\[
  H^1(H,\mathcal O_H^r)
\]
in coherent-sheaf cohomology.  Since \(H\) is affine,
\[
  H^1(H,\mathcal O_H^r)=0
\]
by EGA III, \S1.  Thus \(E'\to H\) has a section.  Consequently,
\(\operatorname{Ext}_{\mathrm{alg}}(H,V)\) is isomorphic to
\(H^2(H,V)\) by Appendix I, \SGARefLink{sga3:i:definition:3-1}{3.1}.  But
\[
  H^i(H,V)=H^i(H,\mathbf W(V))=0\qquad(i>0)
\]
by \SGARefLink{sga3:ix:theorem:3:1}{Exposé IX, 3.1}.  The exact sequence \((**)\) now gives
\[
  \operatorname{Ext}_{\mathrm{alg}}(H,{}_F V)=0,
\]
and hence \(E\) is a trivial extension.

\SGASubsection{5.5}{Proof of \SGARefLink{sga3:xvii:theorem:5:1:1:item:i}{5.1.1(i)}}
\SGARefTarget{sga3:xvii:subsection:5:5}

If \(k\) has characteristic \(0\), then \(U\) is \(k\)-solvable
by \SGARefLink{sga3:xvii:corollary:4:1:3}{4.1.3} and \(H\) is smooth.  It follows from \SGARefLink{sga3:xvii:proposition:5:3:1}{5.3.1} and
\SGARefLink{sga3:xvii:lemma:5:2:3:item:i:d}{5.2.3(d)} that \(E\) is a trivial extension.  The same argument
shows that two lifts of \(H\) in \(E\) are conjugate by an
element of \(U(k)\).

Henceforth we suppose that \(k\) is a field of characteristic
\(p>0\).

\SGAProof{Proof of \SGARefLink{sga3:xvii:lemma:5:5:1:item:i}{(i)}\SGARefLink{sga3:xvii:section:b}{(b)}: the case where \(k\) is perfect and
\(U\) is connected}
We shall reduce to the case in which \(U\) is radicial.  Since
\(k\) is perfect, \(H_{\mathrm{red}}\) is smooth; let \(E'\) be
its inverse image in \(E\).  By \SGARefLink{sga3:xvii:proposition:5:3:1}{5.3.1} and \SGARefLink{sga3:xvii:lemma:5:2:3:item:i}{5.2.3(i)}\SGARefLink{sga3:xvii:appendix:C}{(c)}, the
extension
\[
  1\longrightarrow U\longrightarrow E'
    \longrightarrow H_{\mathrm{red}}\longrightarrow1
\]
is trivial.  Let \(H_1\) be a lift of \(H_{\mathrm{red}}\) in
\(E\).  By Appendix II, \SGARefLink{sga3:xvii:lemma:3-1}{3.1}, there is an integer \(n>0\) such
that
\[
  E^{(n)}=E/{}_{F^n}(E)
\]
is smooth.  Let \(E''\) be the algebraic subgroup of \(E\)
generated by \(H_1\) and \({}_{F^n}(E)\); equivalently, \(E''\)
is the inverse image in \(E\) of the image of \(H_1\) in
\(E^{(n)}\).  Let \(H''\) be the image of \(E''\) in \(H\).  We
claim that \(H''=H\).

Indeed, let \(R\) be the image of \({}_{F^n}(E)\) in \(H\), so
that \(H''\) is generated by \(R\) and \(H_{\mathrm{red}}\).
The group \(H/R\) is a quotient of \(E^{(n)}\), and is therefore
smooth.  Consequently the canonical morphism
\[
  H_{\mathrm{red}}\longrightarrow H/R
\]
is an epimorphism.  Thus \(H\) is generated by \(R\) and
\(H_{\mathrm{red}}\), and hence equals \(H''\).  We have obtained
an exact sequence
\[
  \tag{\(\dagger\)}\label{sga3:xvii:subsection:5-5:equation:dagger}
  1\longrightarrow U''=U\cap E''
    \longrightarrow E''\longrightarrow H\longrightarrow1.
\]
But \(E''\) has the same underlying space as \(H_1\), so \(U''\)
is radicial.  It is plainly enough to prove that the extension
\((\dagger)\) is trivial, and this follows from \SGARefLink{sga3:xvii:proposition:5:4:1}{5.4.1}.

\SGAProof{Proof of \SGARefLink{sga3:xvii:lemma:5:5:1:item:i}{(i)}\SGARefLink{sga3:xvii:section:b}{(b)}: the case where \(k\) is perfect and
\(H\) is connected}
The group \(U\) is an extension of an étale group by its identity
component \(U^0\).  The case where \(U\) is connected having just
been treated, it is enough to consider the case where \(U\) is
étale.  We shall use the following more precise lemma.

\SGAStatement{Lemma 5.5.1}{With the notation of \SGARefLink{sga3:xvii:theorem:5-1-1}{5.1.1}, suppose
in addition that \(U\) is étale.  Then:
\begin{enumerate}
\item[(i)] \SGARefTarget{sga3:xvii:lemma:5:5:1:item:i}
  If \(H\) is connected, it has a unique lift in \(E\), namely
  the identity component \(E^0\) of \(E\).
\item[(ii)] \SGARefTarget{sga3:xvii:lemma:5:5:1:item:ii}
  If \(k\) is algebraically closed, then \(E\) is trivial, and
  any two lifts of \(H\) in \(E\) are conjugate by an element of
  \(U(k)\).
\end{enumerate}}
\SGARefTarget{sga3:xvii:lemma:5:5:1}

\SGAProof{Proof of \SGARefLink{sga3:xvii:lemma:5:5:1:item:i}{(i)}}
Since formation of the identity component commutes with extension
of the base field, we may restrict ourselves to the case where
\(k\) is algebraically closed.  If \(H\) is a torus, then \(E\)
is trivial by \SGARefLink{sga3:xvii:proposition:5:3:1}{5.3.1} and \SGARefLink{sga3:xvii:lemma:5:2:3:item:i}{5.2.3(i)}\SGARefLink{sga3:xvii:section:b}{(b)}, and \(E^0\) is plainly the
unique lift of \(H\).  This already proves that in the general
case \(E^0\cap U\) is radicial.  It is also étale, since \(U\) is
étale, and is therefore zero.  The morphism
\[
  E^0\longrightarrow H
\]
is consequently a monomorphism; it is flat because \(E^0\) is
open in \(E\), and it is surjective because \(H\) is connected.
It is therefore an isomorphism.  If \(H'\) is any other lift of
\(H\), then \(H'\) is connected, hence contained in \(E^0\), and
thus equal to \(E^0\).

\SGAProof{Proof of \SGARefLink{sga3:xvii:lemma:5:5:1:item:ii}{(ii)}}
Let \(H^0\) be the identity component of \(H\).  By \SGARefLink{sga3:xvii:lemma:5:5:1:item:i}{(i)}, \(E^0\)
is the unique lift of \(H^0\) in \(E\).  Put
\[
  E'=E/E^0,\qquad H'=H/H^0.
\]
There is then an extension
\[
  1\longrightarrow U\longrightarrow E'
    \longrightarrow H'\longrightarrow1.
\]
Since \(H'\) is étale, this extension is trivial by
\SGARefLink{sga3:xvii:lemma:5:2:3:item:i}{5.2.3(i)}\SGARefLink{sga3:xvii:section:b}{(b)}.  If \(H'_1\) is a lift of \(H'\) in \(E'\), and
\(H_1\) is its inverse image in \(E\), then \(H_1\) plainly
lifts \(H\) in \(E\).  If \(H_2\) is a second lift of \(H\) in
\(E\), then it contains \(E^0\).  By \SGARefLink{sga3:xvii:lemma:5:2:3:item:i}{5.2.3(i)}\SGARefLink{sga3:xvii:section:b}{(b)}, the image of
\(H_2\) in \(E'\) is conjugate to \(H'_1\) by an element
\(u\in U(k)\), and hence
\[
  H_2=\operatorname{int}(u)H_1.
\]

\SGAStatement{Remark 5.5.2}{Under the hypotheses of \SGARefLink{sga3:xvii:lemma:5:5:1:item:i}{5.5.1(i)},
the group \(E^0\) centralizes \(U\).}
\SGARefTarget{sga3:xvii:remark:5:5:2}

\SGAProof{Proof of \SGARefLink{sga3:xvii:theorem:5:1:1:item:i}{5.1.1(i)}\SGARefLink{sga3:xvii:appendix:A}{(a)}}
Using the composition series
\[
  1\longrightarrow U^0\longrightarrow U
    \longrightarrow U/U^0\longrightarrow1,
\]
assertion \SGARefLink{sga3:xvii:lemma:5:5:1:item:i}{(i)}\SGARefLink{sga3:xvii:appendix:A}{(a)} follows by combining \SGARefLink{sga3:xvii:lemma:5:5:1:item:i}{(i)}\SGARefLink{sga3:xvii:section:b}{(b)} with \SGARefLink{sga3:xvii:lemma:5:5:1:item:ii}{5.5.1(ii)}.

Before proving \SGARefLink{sga3:xvii:theorem:5:1:1:item:i}{5.1.1(i)}\SGARefLink{sga3:xvii:appendix:C}{(c)} and (d), we shall first establish
\SGARefLink{sga3:xvii:theorem:5:1:1:item:ii}{5.1.1(ii)}.

\SGAProof{Proof of \SGARefLink{sga3:xvii:theorem:5:1:1:item:ii}{5.1.1(ii)}\SGARefLink{sga3:xvii:appendix:A}{(a)}}
For want of a satisfactory general statement, we shall describe
a number of cases in which, when \(H\) is smooth, two lifts of
\(H\) in \(E\) are conjugate.

\SGAStatement{\SGAInlineRefTarget{sga3:xvii:proposition:5-6-1}Proposition 5.6.1}{With the notation of \SGARefLink{sga3:xvii:theorem:5-1-1}{5.1.1},
suppose in addition that \(H\) is smooth, and let \(H'\) and
\(H''\) be two lifts of \(H\) in \(E\).  Then \(H'\) and \(H''\)
are conjugate by an element of \(U(k)\) in each of the following
cases:
\begin{enumerate}
\item[(a)] \SGARefTarget{sga3:xvii:proposition:5:6:1:item:a}
  \(k\) is algebraically closed.
\item[(b)] \SGARefTarget{sga3:xvii:proposition:5:6:1:item:b}
  \(U\) is commutative.
\item[(c)] \SGARefTarget{sga3:xvii:proposition:5:6:1:item:c}
  \(k\) is perfect and \(U\) is connected.
\item[(d)] \SGARefTarget{sga3:xvii:proposition:5:6:1:item:d}
  \(U\) is \(k\)-solvable.
\item[(e)] \SGARefTarget{sga3:xvii:proposition:5:6:1:item:e}
  The centralizer of \(H'\) in \(U\) is \(k\)-solvable.
\item[(f)] \SGARefTarget{sga3:xvii:proposition:5:6:1:item:f}
  The group of multiplicative type \(H\) is split by a finite
  Galois extension \(K/k\) of degree prime to \(p\).
\end{enumerate}}
\SGARefTarget{sga3:xvii:proposition:5:6:1}
